% Generated by scripts/assemble.py; corrections belong in source/corrections.json.
\examhead{66}{河北工业大学 2024 年数学分析试卷}{66}
% p067-b05; source PDF page 67
\par\Needspace{5\baselineskip}
\blockanchor{p067-b05}
\majorquestion{一}填空题. 每题 5 分, 共 30 分.\par
% p067-b06; source PDF page 67
\blockanchor{p067-b06}
\question{1}求和函数 $\sum_{n=1}^{\infty} nx^{n-1} = \underline{\quad}$.\par
% p067-b07; source PDF page 67
\blockanchor{p067-b07}
\question{2}求极限  $\lim_{x\to0}\frac{e^{x}\sin x-x(1+x)}{x^{3}}=\underline{\quad}$ .\par
% p067-b08; source PDF page 67
\blockanchor{p067-b08}
\question{3}求球面  $x^{2} + y^{2} + z^{2} = 1$  在点  $(x, y, z) = \left(\frac{\sqrt{2}}{2}, 0, \frac{\sqrt{2}}{2}\right)$  处的切平面方程 \_\_\_\_.\par
% p067-b09; source PDF page 67
\blockanchor{p067-b09}
\question{4}求函数  $f(x)=(x-2)^{\frac{2}{3}}$  在区间  $[-1,2]$  上的最大值和最小值 \_\_\_\_.\par
% p067-b10; source PDF page 67
\blockanchor{p067-b10}
\question{5}设 $S = \{(x, y, z) \mid x^2 + y^2 + z^2 = 1\}$, 方向向外, 求曲面积分 $\iint_{S} x \, \mathrm{d}y \, \mathrm{d}z + y \, \mathrm{d}z \, \mathrm{d}x + z \, \mathrm{d}x \, \mathrm{d}y =$ \_\_\_\_.\par
% p067-b11; source PDF page 67
\blockanchor{p067-b11}
\question{6}求极限  $\lim_{n\to\infty}\left[\frac{1}{n\ln^{2}n}+\frac{1}{(n+1)\ln^{2}(n+1)}+\cdots+\frac{1}{(2n-1)\ln^{2}(2n-1)}\right]=$  \_\_\_\_.\par
% p067-b12; source PDF page 67
\par\Needspace{5\baselineskip}
\blockanchor{p067-b12}
\majorquestion{二}计算题. 每题 10 分, 共 60 分.\par
% p067-b13; source PDF page 67
\blockanchor{p067-b13}
\question{1}判断函数 $f(x) = \frac{x^2}{4} - \cos x$ 的凸性.\par
% p067-b14; source PDF page 67
\par\Needspace{6\baselineskip}
\blockanchor{p067-b14}
\question{2}设 $F$ 为二阶可微的三元函数, 又设由方程\par
\nopagebreak[4]
% p067-b15; source PDF page 67
\blockanchor{p067-b15}
\displayformula{F (x y, y + z, x z) = 0}
% p067-b16; source PDF page 67
\blockanchor{p067-b16}
可以确定隐函数 $z = f(x,y)$ ，求偏导数 $\frac{\partial f}{\partial x}$ 和 $\frac{\partial f}{\partial y}$\par
% p067-b17; source PDF page 67
\blockanchor{p067-b17}
\question{3}设曲面 $\Sigma$ 为球面 $x^{2} + y^{2} + z^{2} = 1$ 在 $z^{2} \geq \frac{1}{2}$ 的部分, 求 $I = \iint_{\Sigma} (1 - z^{2}) \mathrm{d}S$.\par
% p067-b18; source PDF page 67
\par\Needspace{6\baselineskip}
\blockanchor{p067-b18}
\question{4}设 $L$ 是从 $A$ 点 $(2,0)$ 出发, 先沿直线 $x + 2y = 2$ 到 $B$ 点 $(0,1)$, 再沿圆周 $x^{2} + y^{2} = 1$ 到 $C$ 点 $(-1,0)$, 最后由 $C$ 点沿 $x$ 轴到 $A$ 的定向曲线, 计算第二类曲线积分\par
\nopagebreak[4]
% p067-b19; source PDF page 67
\blockanchor{p067-b19}
\displayformula{\int_ {L} (x ^ {2} - 3 y) \mathrm{d} x + (3 x + e ^ {y}) \mathrm{d} y.}
% p067-b20; source PDF page 67
\blockanchor{p067-b20}
\question{5}求解不定积分 $\int \frac{1}{1 + x^3} \mathrm{d}x$.\par
% p067-b21; source PDF page 67
\blockanchor{p067-b21}
\question{6}计算 $\int_0^{+\infty} \frac{e^{-x^2} - e^{-4x^2}}{x^2} \, \mathrm{d}x$.\par
% p067-b22; source PDF page 67
\par\Needspace{5\baselineskip}
\blockanchor{p067-b22}
\majorquestion{三}证明题. 每题15分, 共60分.\par
% p067-b23; source PDF page 67
\blockanchor{p067-b23}
\question{1}设 $f(x) = x^s$，其中参数 $0 < s < 1$，证明：$f(x)$ 在 $[0, +\infty)$ 上一致连续。\par
% p067-b24; source PDF page 67
\par\Needspace{6\baselineskip}
\blockanchor{p067-b24}
\question{2}设 $D$ 为平面上一个有界区域, 它在 $x$ 轴与 $y$ 轴上投影长度分别为 $a$ 和 $b$, $D$ 的面积为 $A$, $(x_0, y_0) \in D$, 证明:\par
\nopagebreak[4]
% p068-b00; source PDF page 68
\blockanchor{p068-b00}
\displayformula{\left| \iint_ {D} (x - x _ {0}) (y - y _ {0}) \mathrm{d} x \mathrm{d} y \right| \leq a b A. \tag {1}}
% p068-b01; source PDF page 68
\blockanchor{p068-b01}
\displayformula{\left| \iint_ {D} (x - x _ {0}) (y - y _ {0})   \mathrm{d} x \mathrm{d} y \right| \leq \frac {1}{4} a ^ {2} b ^ {2}. \tag {2}}
% p068-b02; source PDF page 68
\par\Needspace{6\baselineskip}
\blockanchor{p068-b02}
\question{3}证明函数 $J(x) = \frac{1}{\pi} \int_{0}^{\pi} \cos(t - x \sin(t)) \, \mathrm{d}t$ 满足 Bessel 方程\par
\nopagebreak[4]
% p068-b03; source PDF page 68
\blockanchor{p068-b03}
\displayformula{x ^ {2} J ^ {\prime \prime} (x) + x J ^ {\prime} (x) + (x ^ {2} - 1) J (x) = 0.}
% p068-b04; source PDF page 68
\blockanchor{p068-b04}
\question{4}设函数 $f(x)$ 在 $[a, b]$ 上连续, 证明: $M(x) = \sup_{a \leq t \leq x} f(t)$ 在 $[a, b]$ 上连续.\par
